\documentclass[12pt]{article}

\usepackage[margin=1in]{geometry}
\usepackage{enumitem}
\usepackage{amsmath}
\usepackage{amssymb}
\usepackage{array}
\usepackage{titlesec}

\title{Lecture 9 Notes: Aristotle's \textit{Topics}, Book II}
\author{}
\date{}

\begin{document}

\maketitle

\section*{Central Theme}

Aristotle's \textit{Topics}, Book II begins the practical art of topical
argument. It gives commonplace rules for testing claims about accidental
predication.

\[
\boxed{
\textit{Topics}, \text{ Book II}
=
\text{commonplace rules for arguing about accidental predication}
}
\]

Book I tells us what dialectic is. Book II begins showing us how to do it:
divide the claim, define the terms, test consequences, check opposites, compare
likenesses, and avoid confusing qualified predication with absolute predication.

\section*{1. From Book I to Book II}

Book I introduced:

\begin{itemize}[itemsep=0.25em]
    \item dialectical reasoning,
    \item generally accepted opinions,
    \item propositions and problems,
    \item the four predicables,
    \item the four sources of arguments.
\end{itemize}

Book II now begins the actual use of commonplace rules.

\[
\text{Book I: What dialectic is}
\]

\[
\text{Book II: How to argue about accidents}
\]

\subsection*{Teaching Point}

Book II is where the \textit{Topics} becomes a practical handbook of argument.

\[
\boxed{
\text{Book II begins the practical art of topical argument.}
}
\]

\section*{2. Accident as the First Field of Rules}

The rules in Book II concern accident.

An accident is what may belong or not belong to the same thing.

\[
\boxed{
\text{Accident} = \text{what may belong or not belong to the same subject}
}
\]

Examples:

\[
\text{white}
\]

\[
\text{sitting}
\]

\[
\text{just}
\]

\[
\text{pleasant}
\]

\subsection*{Teaching Point}

Accidental predication is unstable because it may be partial, conditional,
temporary, relative, or qualified.

\[
\boxed{
\text{Accidents are harder to handle than definitions, properties, or genera.}
}
\]

\section*{3. Universal and Particular Problems}

Aristotle begins by distinguishing universal and particular problems.

\subsection*{Universal Problems}

\[
\text{Every pleasure is good.}
\]

\[
\text{No pleasure is good.}
\]

\subsection*{Particular Problems}

\[
\text{Some pleasure is good.}
\]

\[
\text{Some pleasure is not good.}
\]

\subsection*{Teaching Point}

Universal methods also help with particular problems.

If we show that a predicate belongs in every case, then we have also shown that
it belongs in some case.

If we show that it belongs in no case, then we have also shown that it fails in
every case.

\[
\boxed{
\text{Universal proof contains particular proof within it.}
}
\]

\section*{4. Establishing and Overthrowing}

Dialectical argument may aim either to establish a claim or to overthrow it.

\[
\text{Establishing} = \text{showing that the predicate belongs}
\]

\[
\text{Overthrowing} = \text{showing that the predicate does not belong}
\]

Aristotle begins with overthrowing because people commonly assert that
something belongs, while their opponents try to refute the assertion.

\subsection*{Teaching Point}

To refute a universal affirmative claim, one counterexample is enough.

\[
\boxed{
\text{One negative instance overthrows a universal assertion.}
}
\]

\section*{5. Why Conversion Is Precarious in Accidents}

In the case of definition, property, and genus, conversion is more stable.

\[
\text{Definition, property, genus} \rightarrow \text{stable predication}
\]

But accident is different.

\[
\text{Accident} \rightarrow \text{partial, qualified, or conditional predication}
\]

For example, it is not always enough to say:

\[
\text{Justice belongs to this man.}
\]

Someone may reply:

\[
\text{Justice belongs to him only in part.}
\]

\subsection*{Teaching Point}

Accidents can belong in qualified ways.

\[
\boxed{
\text{Accidental predication often requires qualification.}
}
\]

\section*{6. Two Kinds of Error}

Aristotle distinguishes two kinds of error in problems.

\[
\begin{array}{ll}
1. & \text{False statement} \\
2. & \text{Transgression of established diction}
\end{array}
\]

\subsection*{False Statement}

A false statement says that an attribute belongs when it does not.

\[
\text{S is P}
\]

when in fact:

\[
\text{S is not P}
\]

\subsection*{Transgression of Established Diction}

A transgression of diction uses the wrong name for the thing.

Example:

\[
\text{calling a plane-tree a man}
\]

\subsection*{Teaching Point}

Dialectical error can lie either in what is asserted or in how it is named.

\[
\boxed{
\text{Bad predication and bad naming are different failures.}
}
\]

\section*{7. Do Not Mistake Genus for Accident}

A major commonplace rule is to ask whether someone has called an accident what
really belongs in some other way.

This mistake commonly happens with genera.

\subsection*{Example}

It is wrong to say:

\[
\text{White happens to be a colour.}
\]

Colour does not accidentally belong to white.

Rather:

\[
\text{colour is the genus of white.}
\]

\subsection*{Whiteboard}

\[
\text{white} \rightarrow \text{colour}
\]

\[
\text{colour} = \text{genus, not accident}
\]

\subsection*{Teaching Point}

A predicate must be attacked or defended under the right predicable.

\[
\boxed{
\text{First identify whether the predicate is definition, property, genus, or accident.}
}
\]

\section*{8. Examine Cases Species by Species}

When a predicate is asserted or denied universally, Aristotle says to examine
the cases species by species.

Begin with the primary divisions, then proceed down to the more specific cases.

\subsection*{Example}

Suppose someone claims:

\[
\text{The knowledge of opposites is the same.}
\]

Examine kinds of opposition:

\[
\begin{array}{ll}
1. & \text{relative opposites} \\
2. & \text{contraries} \\
3. & \text{privation and possession} \\
4. & \text{contradictories}
\end{array}
\]

Then divide further if necessary.

\[
\text{just and unjust}
\]

\[
\text{double and half}
\]

\[
\text{blindness and sight}
\]

\[
\text{being and not-being}
\]

\subsection*{Teaching Point}

To overthrow a universal claim, find one division where it fails.

\[
\boxed{
\text{Divide the universal until a counterexample appears.}
}
\]

\section*{9. The Constructive Use of Division}

The same rule can also be used constructively.

If the predicate appears to hold in all or many divisions, we may press the
opponent to do one of two things:

\begin{enumerate}[itemsep=0.25em]
    \item assert the universal claim,
    \item or give a negative instance.
\end{enumerate}

\subsection*{Teaching Point}

Division can force the discussion toward either universal admission or
counterexample.

\[
\boxed{
\text{A good division makes the opponent choose between concession and exception.}
}
\]

\section*{10. Define the Accident and the Subject}

Another rule is to define the accident and the subject, then see whether the
definitions expose anything false.

\[
\text{Define the predicate.}
\]

\[
\text{Define the subject.}
\]

\[
\text{Compare definitions.}
\]

\subsection*{Example: Can a God Be Wronged?}

Ask:

\[
\text{What is it to wrong someone?}
\]

Suppose wronging is:

\[
\text{injuring deliberately}
\]

Then a god cannot be wronged, because a god cannot be injured.

\subsection*{Example: Is the Good Man Jealous?}

Ask:

\[
\text{What is jealousy?}
\]

Suppose jealousy is:

\[
\text{pain at the apparent success of a well-behaved person}
\]

Then the good man is not jealous, since that would make him bad.

\subsection*{Teaching Point}

Definitions reveal whether an accidental predication can stand.

\[
\boxed{
\text{When the surface claim is unclear, define the terms.}
}
\]

\section*{11. Continue Defining Until the Issue Is Clear}

Aristotle says we should substitute definitions not only for the original terms,
but also for terms contained in the definitions.

Do not stop until the argument reaches familiar terms.

\subsection*{Teaching Point}

Definitions may need to be unpacked more than once.

\[
\boxed{
\text{A definition that remains obscure has not yet done its dialectical work.}
}
\]

\section*{12. Turn the Problem into a Proposition}

Aristotle also advises us to make the problem into a proposition and then bring
a negative instance against it.

\subsection*{Pattern}

Problem:

\[
\text{Is every } S \text{ P?}
\]

Proposition:

\[
\text{Every } S \text{ is P.}
\]

Attack:

\[
\text{Here is an } S \text{ that is not P.}
\]

\subsection*{Teaching Point}

Dialectic often advances by translating a question into an attackable claim.

\[
\boxed{
\text{Question-form must often become proposition-form before it can be tested.}
}
\]

\section*{13. Ordinary Language and Expert Judgment}

Aristotle distinguishes ordinary naming from expert factual judgment.

We should use names as most people use them.

But when deciding whether something really is of such-and-such a kind, we
should follow those with knowledge.

\subsection*{Example}

We should use the term:

\[
\text{healthy}
\]

as ordinary people use it.

But when deciding whether something tends to produce health, we should listen
to the doctor, not the multitude.

\subsection*{Teaching Point}

Dialectic respects common speech without surrendering factual judgment to the
crowd.

\[
\boxed{
\text{Name by common usage; judge facts by expertise.}
}
\]

\section*{14. Exploit Ambiguity Carefully}

If a term is used in several senses, and the difference is unnoticed, one may
sometimes establish or overthrow a claim by showing it in one sense.

But if the ambiguity is obvious, one must distinguish the senses before arguing.

\subsection*{Example}

Suppose:

\[
\text{right}
\]

means both:

\[
\text{expedient}
\]

and:

\[
\text{honourable}
\]

Then one should try to establish or demolish the claim in both senses. If that
is impossible, one must say clearly which sense is being used.

\subsection*{Teaching Point}

Ambiguity is both a danger and a dialectical opportunity.

\[
\boxed{
\text{Hidden ambiguity may be used; obvious ambiguity must be distinguished.}
}
\]

\section*{15. Distinguish Kinds of ``Of''}

Aristotle analyzes phrases such as:

\[
\text{The science of many things is one.}
\]

The phrase ``of many things'' may mean several things.

\[
\begin{array}{ll}
1. & \text{of an end} \\
2. & \text{of a means to an end} \\
3. & \text{of two ends} \\
4. & \text{of an essential and an accidental attribute}
\end{array}
\]

\subsection*{Examples}

Medicine is of health and dieting.

\[
\text{health} = \text{end}
\]

\[
\text{dieting} = \text{means to the end}
\]

Geometry may be of triangle essentially and of equilateral triangle
accidentally.

\[
\text{triangle} = \text{essential subject}
\]

\[
\text{equilateral} = \text{accidental case of triangle}
\]

\subsection*{Teaching Point}

Relational phrases often hide several logical relations.

\[
\boxed{
\text{The word ``of'' must be interpreted according to the relation involved.}
}
\]

\section*{16. Desire as an Example of Relational Ambiguity}

Aristotle applies the same rule to desire.

\[
\text{desire of X}
\]

may mean:

\[
\text{desire of an end}
\]

\[
\text{desire of a means}
\]

\[
\text{desire of something accidental}
\]

\subsection*{Examples}

\[
\text{desire of health} = \text{desire of an end}
\]

\[
\text{desire of being doctored} = \text{desire of a means}
\]

A sweet-toothed person may desire wine not because it is wine, but because it
is sweet.

\[
\text{desire of wine} = \text{accidental desire}
\]

\subsection*{Teaching Point}

Many relational terms must be divided by the role of their object.

\[
\boxed{
\text{Ask whether the object is desired as end, means, or accident.}
}
\]

\section*{17. Substitute Familiar Terms}

Aristotle says it is useful to replace an obscure or elevated expression with a
more familiar one.

\subsection*{Examples}

\[
\text{exact} \rightarrow \text{clear}
\]

\[
\text{being busy} \rightarrow \text{being fussy}
\]

When the expression becomes more familiar, the thesis becomes easier to attack
or defend.

\subsection*{Teaching Point}

Clarifying language often makes the argument easier to evaluate.

\[
\boxed{
\text{A familiar formulation exposes the force of the claim.}
}
\]

\section*{18. Genus-to-Species Arguments}

Aristotle gives rules for arguing between genus and species.

Some directions of inference are safe; others are not.

\subsection*{Unsafe for Establishing}

From the fact that a genus has an attribute, it does not follow that every
species has that attribute.

\[
\text{Genus has P} \not\Rightarrow \text{Species has P}
\]

Example:

\[
\text{Animal may be flying or quadruped.}
\]

It does not follow that:

\[
\text{Man is flying or quadruped.}
\]

\subsection*{Safe for Overthrowing}

If an attribute does not belong to the genus, it does not belong to the species.

\[
\text{Genus lacks P} \Rightarrow \text{Species lacks P}
\]

\subsection*{Safe from Species to Genus}

If an attribute belongs to the species, it belongs to the genus in some way.

\[
\text{Species has P} \Rightarrow \text{Genus has P}
\]

\subsection*{Teaching Point}

Arguments between genus and species must respect direction.

\[
\boxed{
\text{Do not infer too much from genus to species.}
}
\]

\section*{19. If a Genus Is Predicated, Some Species Must Be Involved}

If a genus is predicated of something, then some species under that genus must
also be involved.

\subsection*{Example: Scientific Knowledge}

If something is scientific knowledge, then it must be some specific science.

\[
\text{scientific knowledge}
\Rightarrow
\text{grammar, music, medicine, geometry, or another science}
\]

\subsection*{Example: Motion}

If the soul is said to move, ask by which species of motion it moves.

\[
\text{motion}
\Rightarrow
\text{generation, destruction, growth, diminution, alteration, or change of place}
\]

If none of these applies, then the soul does not move.

\subsection*{Teaching Point}

Generic claims must be cashed out in specific forms.

\[
\boxed{
\text{If the genus applies, some species must apply.}
}
\]

\section*{20. If Stuck, Attack Definitions}

If one is not well equipped with an argument against the assertion, Aristotle
says to look among definitions, real or apparent.

If one definition is not enough, use several.

\subsection*{Teaching Point}

It is easier to attack people when they are committed to a definition.

\[
\boxed{
\text{Definitions expose the structure of a claim to attack.}
}
\]

\section*{21. Necessary Conditions and Consequences}

Aristotle advises us to ask two questions.

\subsection*{For Establishing}

What condition, if real, would make the thing in question real?

\[
A \Rightarrow B
\]

\[
A
\]

\[
\therefore B
\]

\subsection*{For Overthrowing}

What follows if the thing in question is real?

\[
B \Rightarrow C
\]

\[
\neg C
\]

\[
\therefore \neg B
\]

\subsection*{Teaching Point}

A claim can be tested by its necessary conditions and consequences.

\[
\boxed{
\text{To establish, find what entails it; to overthrow, refute what it entails.}
}
\]

\section*{22. Check the Time Involved}

Aristotle says to look at the time involved in the predication.

\subsection*{Example: Nourishment and Growth}

Someone may say:

\[
\text{What is being nourished necessarily grows.}
\]

But animals are always being nourished, while they do not always grow.

\subsection*{Example: Knowing and Remembering}

Someone may say:

\[
\text{Knowing is remembering.}
\]

But remembering concerns the past, while knowing can concern:

\[
\text{past, present, and future}
\]

\subsection*{Teaching Point}

Temporal mismatch can expose a false predication.

\[
\boxed{
\text{Always ask whether the predicate and subject belong to the same temporal range.}
}
\]

\section*{23. Diverting the Argument}

Aristotle discusses a sophistical turn in which the questioner directs the
argument toward another statement against which he is well supplied with
arguments.

This may be:

\[
\begin{array}{ll}
1. & \text{really necessary} \\
2. & \text{apparently necessary} \\
3. & \text{neither really nor apparently necessary}
\end{array}
\]

\subsection*{Teaching Point}

A dialectician must distinguish relevant redirection from irrelevant side issue.

\[
\boxed{
\text{Not every victory on a side issue is a victory over the thesis.}
}
\]

\section*{24. Use Consequences of a Statement}

Anyone who makes one statement has, in a sense, made several statements,
because the statement has necessary consequences.

\subsection*{Example}

If someone says:

\[
X \text{ is a man}
\]

then he has also implied that:

\[
X \text{ is an animal}
\]

\[
X \text{ is animate}
\]

\[
X \text{ is biped}
\]

\[
X \text{ is capable of reason and knowledge}
\]

\subsection*{Teaching Point}

To attack a claim, attack what it necessarily implies.

\[
\boxed{
\text{Demolishing a necessary consequence demolishes the original claim.}
}
\]

\section*{25. Use Exhaustive Alternatives}

If a subject must have one and only one of two predicates, then proving or
disproving one gives material for the other.

\subsection*{Example}

A human body may have:

\[
\text{health}
\]

or:

\[
\text{disease}
\]

If one is present, the other is absent.

If one is absent, the other is present.

\[
S \text{ must be either } P \text{ or } Q
\]

\[
P \Rightarrow \neg Q
\]

\[
\neg P \Rightarrow Q
\]

\subsection*{Teaching Point}

Where alternatives are exhaustive, attacking one side supports the other.

\[
\boxed{
\text{Exclusive alternatives make arguments convertible.}
}
\]

\section*{26. Reinterpret Terms Literally}

Aristotle notes that one may devise an attack by taking a term in its literal
meaning rather than its established meaning.

\subsection*{Examples}

\[
\text{strong at heart}
\]

may be taken literally as:

\[
\text{having a strong heart}
\]

instead of:

\[
\text{being courageous}
\]

Likewise:

\[
\text{well-starred}
\]

may be interpreted as:

\[
\text{having a good star}
\]

\subsection*{Teaching Point}

Literal reinterpretation can expose tensions in language, though it may also
become merely verbal.

\[
\boxed{
\text{Dialectic must distinguish real analysis from verbal play.}
}
\]

\section*{27. Necessary, Usual, and Chance}

Aristotle distinguishes three modes of occurrence.

\[
\begin{array}{ll}
1. & \text{Necessary} \\
2. & \text{Usual} \\
3. & \text{Chance}
\end{array}
\]

Some things occur of necessity.

Some things occur usually.

Some things occur however it may chance.

\subsection*{Rules}

If a necessary event is said to occur merely usually, that is a mistake.

If a usual event is said to occur necessarily, that is also a mistake.

If a chance event is said to occur necessarily or usually, that is a mistake.

\subsection*{Whiteboard}

\[
\text{necessary} \neq \text{usual} \neq \text{chance}
\]

\subsection*{Teaching Point}

Many errors come from assigning the wrong modal strength to a predicate.

\[
\boxed{
\text{Ask whether the claim is necessary, usual, or merely occasional.}
}
\]

\section*{28. Do Not Predicate an Accident of Itself}

Aristotle warns against treating the same thing as different merely because it
has different names.

\subsection*{Example}

Prodicus divided pleasures into:

\[
\text{joy}
\]

\[
\text{delight}
\]

\[
\text{good cheer}
\]

If these are merely names of the same thing, then it would be a mistake to say:

\[
\text{Joyfulness is an accident of cheerfulness.}
\]

That would be making something an accident of itself.

\subsection*{Teaching Point}

Different names do not necessarily mean different things.

\[
\boxed{
\text{Do not mistake verbal plurality for real distinction.}
}
\]

\section*{29. Rules from Contraries}

Aristotle next turns to contraries.

Contrary verbs and contrary objects can be joined in several ways.

\subsection*{Examples}

\[
\text{doing good to friends}
\]

\[
\text{doing evil to enemies}
\]

\[
\text{doing evil to friends}
\]

\[
\text{doing good to enemies}
\]

Some of these are contrary to each other, and some are not.

\subsection*{Teaching Point}

The same course may have more than one contrary.

\[
\boxed{
\text{Select whichever contrary is useful in attacking the thesis.}
}
\]

\section*{30. If an Accident Has a Contrary, Test the Contrary}

If the accident asserted has a contrary, see whether the contrary belongs to
the same subject.

If the contrary belongs, then the original accident cannot belong.

\subsection*{Pattern}

\[
S \text{ has } P
\]

\[
P \text{ is contrary to } Q
\]

\[
S \text{ has } Q
\]

\[
\therefore S \text{ does not have } P
\]

\subsection*{Teaching Point}

Contrary predicates cannot belong at the same time to the same thing.

\[
\boxed{
\text{Contraries test whether an accident can belong.}
}
\]

\section*{31. Does the Subject Admit the Contrary?}

If an asserted accident has a contrary, ask whether the subject can admit the
contrary.

This is useful especially for establishing possibility.

\subsection*{Example}

If the subject can admit friendship, then it may also be capable of admitting
hatred, depending on the kind of subject involved.

\subsection*{Teaching Point}

Showing that the subject can admit the contrary may show that the asserted
accident is possible, even if not actual.

\[
\boxed{
\text{Capacity for contraries helps test possible predication.}
}
\]

\section*{32. The Four Modes of Opposition}

Aristotle uses four modes of opposition.

\[
\begin{array}{ll}
1. & \text{Contradiction} \\
2. & \text{Contrariety} \\
3. & \text{Privation and possession} \\
4. & \text{Relative opposition}
\end{array}
\]

\subsection*{Teaching Point}

Opposition supplies structured ways to reverse or test a claim.

\[
\boxed{
\text{Different oppositions yield different argumentative routes.}
}
\]

\section*{33. Arguments from Contradictories}

In contradiction, the sequence is converted.

\subsection*{Example}

If:

\[
\text{man} \Rightarrow \text{animal}
\]

then:

\[
\text{not animal} \Rightarrow \text{not man}
\]

\subsection*{Pattern}

\[
S \Rightarrow P
\]

\[
\neg P \Rightarrow \neg S
\]

\subsection*{Teaching Point}

Contradictory opposition gives a powerful form of reversal.

\[
\boxed{
\text{If the predicate is denied, the subject is denied.}
}
\]

\section*{34. Arguments from Contraries}

In contraries, the sequence is often direct.

\subsection*{Example}

\[
\text{courage} \rightarrow \text{virtue}
\]

\[
\text{cowardice} \rightarrow \text{vice}
\]

Likewise:

\[
\text{courage} \rightarrow \text{desirable}
\]

\[
\text{cowardice} \rightarrow \text{objectionable}
\]

\subsection*{Teaching Point}

If the contrary of one term follows the contrary of another, that may support
or undermine the original claim.

\[
\boxed{
\text{Contrary sequences help test the original sequence.}
}
\]

\section*{35. Privation and Possession}

Privation and possession also yield direct sequences.

\subsection*{Example}

\[
\text{sight} \rightarrow \text{sensation}
\]

\[
\text{blindness} \rightarrow \text{absence of sensation}
\]

Sight is a possession; blindness is a privation.

\subsection*{Teaching Point}

Privation and possession are not mere contradiction. They involve the absence
of a state naturally apt to be present.

\[
\boxed{
\text{Privation arguments track the loss or absence of a proper state.}
}
\]

\section*{36. Relative Opposition}

Relative terms should also be studied.

\subsection*{Example}

If:

\[
3/1 \text{ is a multiple}
\]

then:

\[
1/3 \text{ is a fraction}
\]

Similarly:

\[
\text{knowledge} \rightarrow \text{object of knowledge}
\]

\[
\text{sight} \rightarrow \text{object of sight}
\]

\subsection*{Teaching Point}

Relative terms bring their correlates with them.

\[
\boxed{
\text{To assert one relative is often to imply its corresponding relative.}
}
\]

\section*{37. Co-Ordinates and Inflected Forms}

Aristotle says to examine co-ordinates and inflected forms.

\subsection*{Co-Ordinates}

Co-ordinates are terms in the same kindred series.

\[
\text{justice}
\]

\[
\text{just deeds}
\]

\[
\text{just man}
\]

\subsection*{Inflected Forms}

Inflected forms include adverbial or related expressions.

\[
\text{justly}
\]

\[
\text{courageously}
\]

\[
\text{healthily}
\]

\subsection*{Teaching Point}

A predicate often travels through a family of related terms.

\[
\boxed{
\text{If one member of a term-family is praiseworthy, others may be also.}
}
\]

\section*{38. Test the Contrary Co-Ordinates}

Aristotle says to look not only at the co-ordinates of the subject mentioned,
but also at the co-ordinates of its contrary.

\subsection*{Example}

If:

\[
\text{justice is knowledge}
\]

then one may test whether:

\[
\text{injustice is ignorance}
\]

If:

\[
\text{justly means knowingly}
\]

then test whether:

\[
\text{unjustly means ignorantly}
\]

\subsection*{Teaching Point}

A proposed relation among terms should survive comparison with the contrary
term-family.

\[
\boxed{
\text{Term-families can confirm or expose a predication.}
}
\]

\section*{39. Generation and Destruction}

Aristotle advises looking at how something comes to be and passes away.

\subsection*{Generation}

If the mode of generation is good, the thing generated may be good.

\[
\text{good generation} \Rightarrow \text{good thing}
\]

If the mode of generation is evil, the thing generated may be evil.

\[
\text{evil generation} \Rightarrow \text{evil thing}
\]

\subsection*{Destruction}

In destruction, the relation is reversed.

If the destruction of something is good, then the thing destroyed may be evil.

\[
\text{good destruction} \Rightarrow \text{evil thing destroyed}
\]

If the destruction of something is evil, then the thing destroyed may be good.

\[
\text{evil destruction} \Rightarrow \text{good thing destroyed}
\]

\subsection*{Teaching Point}

How something comes to be or passes away can reveal its value.

\[
\boxed{
\text{Generation and destruction provide signs of goodness and badness.}
}
\]

\section*{40. Productive and Destructive Causes}

The same reasoning applies to causes that produce or destroy.

\[
\text{good productive cause} \Rightarrow \text{good product}
\]

\[
\text{good destructive cause} \Rightarrow \text{bad thing destroyed}
\]

\subsection*{Teaching Point}

Causes can be used dialectically to test the value of their effects.

\[
\boxed{
\text{Examine what produces or destroys the thing in question.}
}
\]

\section*{41. Arguments from Likeness}

Book II continues the use of likeness from Book I.

If one similar case has an attribute, the other may as well.

\subsection*{Example}

If one branch of knowledge has more than one object, perhaps opinion does too.

If possessing sight is seeing, then possessing hearing may be hearing.

\[
\text{As } A \text{ is to } B, \text{ so } C \text{ may be to } D.
\]

\subsection*{Teaching Point}

Analogical likeness gives arguments, but analogy must be tested.

\[
\boxed{
\text{Similarity supplies a route of argument, not an automatic proof.}
}
\]

\section*{42. Single and Multiple Cases}

Aristotle warns that single and multiple cases may differ.

\subsection*{Example}

If someone says:

\[
\text{to know a thing is to think of it}
\]

then one might infer:

\[
\text{to know many things is to think of many things}
\]

But this is false.

A person may know many things without actively thinking of all of them.

\subsection*{Teaching Point}

An analogy between one and many must be tested.

\[
\boxed{
\text{What holds for one case may not hold for many cases together.}
}
\]

\section*{43. Arguments from Greater and Less}

Aristotle gives rules from greater and less degrees.

\subsection*{Rule 1}

See whether a greater degree of the predicate follows a greater degree of the
subject.

\[
\text{greater pleasure} \rightarrow \text{greater good}
\]

\[
\text{greater wrong} \rightarrow \text{greater evil}
\]

\subsection*{Teaching Point}

If increase in the subject carries increase in the predicate, that supports the
predication.

\[
\boxed{
\text{Greater subject, greater predicate: this supports belonging.}
}
\]

\section*{44. More Likely and Less Likely Cases}

Another rule:

If a predicate does not belong where it is more likely to belong, then it does
not belong where it is less likely.

\[
\text{not where more likely}
\Rightarrow
\text{not where less likely}
\]

If it belongs where it is less likely, then it belongs where it is more likely.

\[
\text{yes where less likely}
\Rightarrow
\text{yes where more likely}
\]

\subsection*{Teaching Point}

Greater-and-less arguments exploit comparative plausibility.

\[
\boxed{
\text{The less likely case can prove the more likely case; the more likely failure can refute the less likely.}
}
\]

\section*{45. Like Degrees}

The same kind of argument applies to like degrees.

If a predicate belongs to one subject in a like degree, it may belong to the
other as well.

If it does not belong to one, it may fail to belong to the other as well.

\subsection*{Teaching Point}

Where cases are equal in likelihood, one case can be used as a test for the
other.

\[
\boxed{
\text{Like degree supports like predication.}
}
\]

\section*{46. Addition and Intensification}

Aristotle says we may argue from the addition of one thing to another.

If adding something makes a whole good or white, then the added thing may
itself be good or white.

\[
X + Y \rightarrow \text{good whole}
\]

\[
\therefore Y \text{ may be good}
\]

If adding something intensifies a character already present, the thing added
has that character too.

\subsection*{Teaching Point}

Addition can reveal the character of what is added.

\[
\boxed{
\text{What intensifies a predicate may possess that predicate.}
}
\]

\section*{47. The Limit of the Addition Rule}

The addition rule is not convertible for overthrowing.

If adding a good thing to an evil thing does not make the whole good, it does
not follow that the added thing is not good.

\[
\text{good} + \text{evil} \not\Rightarrow \text{good whole}
\]

Likewise:

\[
\text{white} + \text{black} \not\Rightarrow \text{white whole}
\]

\subsection*{Teaching Point}

Failure to transform a whole does not disprove the character of the added
part.

\[
\boxed{
\text{A good part may fail to make a bad whole good.}
}
\]

\section*{48. Greater and Less Imply Absolute Predication}

Any predicate of which we can speak in greater and less degrees belongs
absolutely.

If something is more white, then it is white.

If something is more good, then it is good.

\[
\text{more P} \Rightarrow \text{P}
\]

\subsection*{Teaching Point}

Degree implies possession.

\[
\boxed{
\text{Nothing can be more P unless it is P.}
}
\]

\section*{49. Absolute Predication Does Not Always Admit Degree}

The converse does not hold.

Some predicates belong absolutely but do not admit greater and less.

\subsection*{Example}

\[
\text{man}
\]

A man is a man, but one man is not more man than another in the relevant sense.

\subsection*{Teaching Point}

Possession does not always imply degree.

\[
\boxed{
\text{P} \not\Rightarrow \text{more-or-less P}
}
\]

\section*{50. Qualified and Absolute Predication}

The end of Book II distinguishes what is said absolutely from what is said with
qualification.

Something may be true:

\[
\text{in a respect}
\]

\[
\text{at a time}
\]

\[
\text{in a place}
\]

\[
\text{under a condition}
\]

without being true absolutely.

\subsection*{Examples}

It may be good to take medicine when ill.

\[
\text{good when ill} \neq \text{good absolutely}
\]

It may be honorable among the Triballi to sacrifice one's father.

\[
\text{honorable among the Triballi} \neq \text{honorable absolutely}
\]

\subsection*{Teaching Point}

Do not confuse qualified truth with absolute truth.

\[
\boxed{
\text{Qualified predication must not be treated as absolute predication.}
}
\]

\section*{51. What Is Said Absolutely}

A thing is said absolutely when one is prepared to say it without adding any
qualification.

\subsection*{Example}

One may say without addition:

\[
\text{To honour the gods is honourable.}
\]

But one should not say without addition:

\[
\text{To sacrifice one's father is honourable.}
\]

If it is honourable only among certain people or under certain conditions, then
it is not honourable absolutely.

\subsection*{Teaching Point}

Absolute predication is predication without addition.

\[
\boxed{
\text{What is true absolutely needs no qualifying phrase.}
}
\]

\section*{52. The Architecture of Book II}

Book II moves through a series of practical rules.

\[
\text{universal and particular problems}
\rightarrow
\text{why accident is unstable}
\rightarrow
\text{errors and wrong predicables}
\rightarrow
\text{division and counterexample}
\]

\[
\rightarrow
\text{definitions and ambiguity}
\rightarrow
\text{genus/species rules}
\rightarrow
\text{conditions and consequences}
\rightarrow
\text{time and modality}
\]

\[
\rightarrow
\text{opposition}
\rightarrow
\text{co-ordinates, generation, likeness}
\rightarrow
\text{greater and less}
\rightarrow
\text{qualified vs. absolute}
\]

\subsection*{Teaching Point}

Book II is Aristotle's toolkit for testing accidental predications.

\[
\boxed{
\text{Book II teaches how to handle predicates that may belong only conditionally.}
}
\]

\section*{53. Summary of Main Ideas}

\begin{enumerate}[itemsep=0.5em]
    \item Book II begins the practical commonplace rules of the
    \textit{Topics}.

    \item The main concern is accidental predication.

    \item Universal and particular problems are connected because universal
    proof includes particular proof.

    \item Accidental predication is difficult because it may be partial,
    qualified, temporary, or conditional.

    \item Dialectical errors may arise from false assertion or from misuse of
    established diction.

    \item A predicate must be classified under the right predicable.

    \item Universal claims should be tested species by species.

    \item Definitions can expose whether an accidental predication is possible.

    \item Ambiguous terms must be distinguished when their ambiguity matters.

    \item Relational terms such as ``of'' may hide several different logical
    relations.

    \item Arguments from genus to species must respect direction.

    \item If a genus applies, some species under that genus must apply.

    \item Claims can be attacked through their necessary consequences.

    \item Temporal mismatch can refute a proposed predication.

    \item Necessary, usual, and chance predications must be distinguished.

    \item Different names do not always indicate different things.

    \item Contraries, contradictories, privation, and relatives provide
    structured lines of attack.

    \item Co-ordinates and inflected forms help test related predications.

    \item Generation, destruction, and causes can reveal value.

    \item Likeness provides arguments by analogy, but analogy must be tested.

    \item Greater-and-less reasoning works through comparative plausibility.

    \item Qualified truth must not be confused with absolute truth.
\end{enumerate}

\section*{Final Framing}

\[
\boxed{
\textit{Topics}, \text{ Book II, is Aristotle's toolkit for testing accidental predications.}
}
\]

Book I tells us what dialectic is. Book II begins showing us how to do it:
divide the claim, define the terms, test the consequences, check the opposites,
compare likenesses, and never confuse the qualified with the absolute.

\[
\boxed{
\text{Accidental claims are tested by qualification, division, opposition, and consequence.}
}
\]

\end{document}